\documentclass[11pt]{article}
\usepackage[T1]{fontenc}
\usepackage{lmodern}
\usepackage[margin=1in]{geometry}
\usepackage{amsmath,amssymb,amsthm}
\usepackage{enumitem}
\usepackage[colorlinks=true,allcolors=blue]{hyperref}
\newtheorem{theorem}{Theorem}
\newtheorem{lemma}[theorem]{Lemma}
\newtheorem{corollary}[theorem]{Corollary}
\newcommand{\R}{\mathbb R}
\newcommand{\Q}{\mathbb Q}
\newcommand{\LAS}{\operatorname{LAS}}
\newcommand{\GAS}{\operatorname{GAS}}
\newcommand{\HGAS}{\operatorname{HGAS}}
\title{Local asymptotic stability of rational polynomial ODEs is undecidable}
\author{Kaiwen Jin\\\small\href{mailto:yui.ui.siki@gmail.com}{yui.ui.siki@gmail.com}}
\date{October 12, 2026}
\hypersetup{pdftitle={Local asymptotic stability of rational polynomial ODEs is undecidable},pdfauthor={Kaiwen Jin},pdfsubject={A literature resolution of AIM Problem 338}}

\begin{document}
\maketitle

\begin{abstract}
No Turing algorithm can decide local asymptotic stability of the origin for
every exactly encoded rational polynomial autonomous ODE, when dimension and
degree are unrestricted. This follows from the recent global
asymptotic-stability theorem of Shevchenko, Liu and Liu
\cite{SLL}. We give the reduction for the precise formulation of AIM
Problem~338, including forward existence near the equilibrium. For a
homogeneous polynomial vector field, local asymptotic stability implies global
asymptotic stability by rescaling both the initial condition and time.
\end{abstract}

\section{The question and the cited theorem}

The question is whether exact polynomial coefficients suffice for a terminating
decision procedure for local asymptotic stability. The distinction between
stability and attraction matters: a stable equilibrium whose nearby solutions
do not converge to it is a negative instance. AIM Problem~338
\cite{AIM338} asks this question for rational polynomial continuous-time
systems without a bound on dimension or degree.

Shevchenko, Liu and Liu \cite{SLL} prove undecidability of global asymptotic
stability even for a suitable fixed dimension and fixed homogeneous degree.
This note records the consequence for the full local question. We use
Theorem~1.5 of their October~8, 2026 preprint, version~1; the undecidability
construction is theirs. The additional argument here is the elementary
local-to-global implication for homogeneous fields and the exact conversion
of the decision problems.

An input consists of an integer $n\geq1$ and a finite monomial description of
\[
  f\in\Q[x_1,\ldots,x_n]^n,\qquad f(0)=0,
  \qquad \dot x=f(x)\quad\text{on }\R^n.
\]
Coefficients are exact rationals and exponents are nonnegative integers.
Write $\phi_f(t,x_0)$ for the unique maximal classical solution from $x_0$,
and let $\|\cdot\|$ be the Euclidean norm. We use precisely the following
meaning of $\LAS(f)$:
\begin{enumerate}[label=(\roman*),leftmargin=*]
\item For every $\varepsilon>0$ there is a $\delta>0$ such that, if
$\|x_0\|<\delta$, the solution exists for all $t\geq0$ and
$\|\phi_f(t,x_0)\|<\varepsilon$ for all $t\geq0$.
\item There is an $r>0$ such that, if $\|x_0\|<r$, the solution exists for
all $t\geq0$ and $\phi_f(t,x_0)\to0$ as $t\to\infty$.
\end{enumerate}
Thus forward completeness is required near zero; it is not assumed for
arbitrary initial states. A decider must halt on every valid finite input and
return the truth value of this conjunction.

The predicate $\GAS(f)$ requires $f(0)=0$, condition~(i), and existence for all
$t\geq0$ with convergence to zero for every $x_0\in\R^n$. This agrees with
the convention in \cite[\S1.1]{SLL}.

For fixed positive integers $N,d$, let $\HGAS(N,d)$ be the set of rational
coefficient codes for homogeneous degree-$d$ polynomial vector fields on
$\R^N$ satisfying $\GAS$. Invalid codes are excluded. A set is recursively
enumerable if a Turing machine halts precisely on its members. It is
many-one complete if it is recursively enumerable and, for each recursively
enumerable set $A$, there is a total computable function $q$ such that
$e\in A$ if and only if $q(e)$ belongs to that set.

\begin{theorem}[{Shevchenko--Liu--Liu \cite[Theorem~1.5, p.~4]{SLL}}]
There are a fixed integer $N_*\geq1$ and a fixed odd integer $d_*\geq3$ for
which $\HGAS(N_*,d_*)$ is many-one complete.
\end{theorem}

The coefficients and the autonomous continuous-time model in this theorem
are the same as those above. Its proof is in \cite[\S15, pp.~44--45]{SLL}.

\section{From local to global stability}

Homogeneity allows any initial state to be moved into the local attraction
ball. The corresponding change of time preserves the entire forward half-line.

\begin{lemma}\label{lem:scaling}
Let $f:\R^n\to\R^n$ be a homogeneous polynomial vector field of degree
$d\geq1$. Then $\LAS(f)$ holds if and only if $\GAS(f)$ holds, with the
definitions above.
\end{lemma}

\begin{proof}
Global asymptotic stability implies both local conditions, with any positive
attraction radius. Conversely, assume $\LAS(f)$ and let $r>0$ be supplied by
condition~(ii). Condition~(i) already gives the stability required for $\GAS$.

Fix $x_0\ne0$ and choose
\[
  \lambda=\min\left\{1,\frac{r}{2\|x_0\|}\right\}>0.
\]
Then $\|\lambda x_0\|\leq r/2<r$, so
$y(t)=\phi_f(t,\lambda x_0)$ exists for every $t\geq0$ and tends to zero.
Define, on the whole half-line,
\[
  x(t)=\lambda^{-1}y(\lambda^{1-d}t).
\]
The chain rule and $f(\lambda x)=\lambda^d f(x)$ give
\[
  x'(t)=\lambda^{-d}f\bigl(y(\lambda^{1-d}t)\bigr)
       =\lambda^{-d}f\bigl(\lambda x(t)\bigr)=f(x(t)),
  \qquad x(0)=x_0.
\]
Since $\lambda^{1-d}$ is a fixed positive number, $x$ is defined for all
nonnegative times and $x(t)\to0$. By uniqueness it agrees with the maximal
solution from $x_0$ on their common domain. This global continuation rules
out a finite forward endpoint of that maximal solution.

For $x_0=0$ the solution is identically zero, since $f(0)=0$. Hence every
initial state has a forward-complete solution converging to zero. Together
with condition~(i), this proves $\GAS(f)$.
\end{proof}

\section{The decision problem}

\begin{corollary}\label{cor:aim}
There is no Turing algorithm that halts on every finite exact rational
polynomial input of unrestricted dimension and degree with $f(0)=0$ and
correctly decides conditions~(i) and~(ii). In particular, the answer to AIM
Problem~338 at the revision in \cite{AIM338} is negative.
\end{corollary}

\begin{proof}
Suppose such a decider $D$ exists. Fix $N_*,d_*$ from Theorem~1. We construct
a total decider $C$ for $\HGAS(N_*,d_*)$.

On an arbitrary source code, first check the finite rational coefficient
syntax and the required dimension and degree. Combine any repeated monomials
by exact rational addition, then verify that each nonzero monomial has total
degree $d_*$. Zero coordinate polynomials are allowed. The checks terminate;
on an invalid code, return false, as required by the definition of $\HGAS$.
In a format that stores only coefficients, the fixed dimension and degree
are hardcoded. If they are explicit metadata, check them as well.

On a valid code, write the same vector field as the finite monomial lists
accepted by $D$. A dense degree-$d_*$ coefficient list can be converted by
enumerating its finitely many monomials and omitting zero coefficients. No
coefficient is approximated. As $d_*\geq1$, the resulting field vanishes at
zero. Thus the input to $D$ is valid and $D$ halts. Lemma~\ref{lem:scaling}
identifies its answer about $\LAS$ with the required answer about $\GAS$.
The resulting $C$ is total. The numerical values of the fixed existential
constants need not be extracted: their existence suffices to give a finite
program with those constants hardcoded.

Let $K$ be the set of codes of pairs $(M,w)$ for which the Turing machine $M$
halts on $w$. It is recursively enumerable by simulation. It is undecidable:
given a hypothetical total halting decider $H$, a machine $B$ on a machine
code $u$ could loop if $H(u,u)$ answers yes and halt if it answers no. On its
own code $b$, either answer $H(b,b)$ would then be wrong.

By Theorem~1, a total computable map $q$ reduces $K$ to
$\HGAS(N_*,d_*)$. Running $q$ and then $C$ would therefore decide $K$ on
every input, a contradiction.
\end{proof}

\paragraph{AI assistance.}
OpenAI Codex (gpt-6.1-sol) assisted with literature retrieval, drafting and
checking this note. The cited undecidability theorem is due to the authors
of \cite{SLL}.

\begin{thebibliography}{9}
\bibitem{SLL}
I.~O. Shevchenko, J.~Liu and X.~Liu,
\emph{Global asymptotic stability of homogeneous polynomial vector fields is
undecidable}, preprint, arXiv:2610.12434v1 [math.DS], October~8, 2026.
\href{https://arxiv.org/abs/2610.12434v1}{arXiv:2610.12434v1}.

\bibitem{AIM338}
AIM problem catalogue,
\emph{338: Decidability of local asymptotic stability for polynomial ODEs},
revision \texttt{c9929805cece044705abaaebe1cebddac1ba11a6}.
\href{https://github.com/MColbrook/AIM/blob/c9929805cece044705abaaebe1cebddac1ba11a6/problems/338-polynomial-local-stability-decidability.md}{Immutable problem statement}.
\end{thebibliography}
\end{document}
